% titel: Blatt 9, A2: Glycerin-Lösung, Gefrier-/Siedepunkt
% Teil II: der durchgerechnete Weg.

\subsection*{Gegeben und gesucht}
2.000 g Glycerin (C$_3$H$_8$O$_3$) in 120.0 g Wasser. Gesucht $T_f$ und
$T_b$. $E_f=1.86$, $E_b=0.513\ K\,kg\,mol^{-1}$.

\subsection*{Rechnung}
$M(\text{Glycerin}) = 3\cdot12.01+8\cdot1.008+3\cdot16.00 \approx 92.09$
g/mol.
\[
  n = \frac{2.000}{92.09} = 0.02172\ mol,\qquad
  b = \frac{0.02172}{0.120} = 0.1810\ mol\,kg^{-1}.
\]
(Nicht-Elektrolyt, $i=1$.)
\begin{align*}
  \Delta T_f &= 1\cdot1.86\cdot0.1810 \approx 0.337\ K,\\
  \Delta T_b &= 1\cdot0.513\cdot0.1810 \approx 0.0928\ K.
\end{align*}

\begin{ergebnis}
\[
  T_f = 0 - 0.337 \approx -0.34\,°C,\qquad
  T_b = 100 + 0.093 \approx 100.09\,°C
  \quad(\text{bei Normaldruck}).
\]
\end{ergebnis}

\begin{kontrolle}
Die Erniedrigung (0.34 K) wirkt deutlich, die Siedepunkterhöhung
(0.09 K) ist klein --- passend, denn $E_f > E_b$. Vorzeichen: beim
Gefrieren abziehen, beim Sieden addieren.
\end{kontrolle}

\begin{merksatz}
$\Delta T = E\cdot b$ mit $b$ in mol/kg Lösungsmittel. Gefrierpunkt sinkt,
Siedepunkt steigt. Nicht-Elektrolyte: $i=1$.
\end{merksatz}
